\section{Corrections recorded}\label{app:corrections}

Corrections are part of the record of any audit. This appendix lists them neutrally: first those to the audit's own notes and to the two Claude instances' earlier statements, then those made in the programme's files themselves, then the corrections of printed sources, then the corrections and generalisations of version~2 of this reader, and finally the changes of version~3. None of the corrections in the first three subsections changes a theorem stated in Sections~\ref{sec:sheets}--\ref{sec:other}; where a correction changed the scope of a statement, the statement in this reader is the corrected one.

\subsection{Corrections to the audit's own statements}

\begin{itemize}
\item \emph{An earlier claim by claude-ab was false.} A note of 23 September claimed that zeros of $\zeta(s,1+t)$ lie on vertical lines only at the Eulerian times. Certified crossings computed in the programme refuted it: zero~4 of $\zeta$ crosses $\re s=\tfrac12$ at $t=0.95095300425$ and zero~5 at $t=0.858704259442$ (located numerically at 25 digits). The erratum, \note{ERRATUM\_FLIP\_FABLE\_ADDENDUM\_4.md}, is among the withdrawn notes in the branch history cited in the introduction, and its check is \note{checks/zeta/hurwitz\_crossing\_check.py}.
\item \emph{A table of claude-ab} labelled as the first 2000 zeros at $T=2515$ contains 1999 zeros; the entries themselves are right for $T=2515$.
\item \emph{A zero count} in \note{08\_} was 13 where it should have been 24; found by a referee pass and by the programme's own reading, and corrected together with three citation fixes.
\item \emph{Referee passes.} Nineteen passes over the audit notes and a two-part final pass over the summaries were made by Claude instances that had not written the notes. They found, among many wording and attribution points, the following statements false or overstated as written, all corrected before this reader was written: a Hurwitz flow described as translating the zero set (only its first jet vanishes on $\Zs-1$); a pair of objects described as the two copies of the zero set when one of them holds only off-line data; a lemma false without the hypothesis $\delta\le1$ (the Gram sandwich); a heuristic about the Yang--Mills mass gap stated as a fact; the term \emph{determinant weight} used for the average weight of a whole quartet, whereas Weil~II's determinant weights belong to the constituents; the $\ZZ_4$ anomaly said to vanish identically (it equals the one-line parity); Turing's method said to give the bulk count exactly (Turing-type data give $P(T)+\sum\lfloor m/2\rfloor$); the first Hurwitz jet's value at $s=0$ (it is $-1$); and a numerical observation at $a=2$ stated as a theorem. Each note lists the findings applied to it, and \note{10\_RESULTS\_DIGEST\_DRAFT.md} §6 lists them pass by pass.
\item \emph{Priority.} The general formal-logarithm criterion and the sharp bound $-1/6$ were derived first in the programme (\lbl{CEI} Theorem 3.1 and \lbl{ETR1}); claude-ab's derivation is later and independent (\note{29\_} §1). The priority of the programme's \lbl{NPE4.2} over part of \note{31\_} Lemma 31.2 is likewise recorded.
\end{itemize}

\subsection{Corrections in the programme's files}

\begin{itemize}
\item \emph{The programme's own late corrections} (final reports of its working sessions of 24--25 September) are carried by the versions audited here: the real boundary pole of the tensor Euler functions gives the position and dimension data only for conjugation-stable exponent data in even degrees (\lbl{ETR9}); zero multiplicity gives vanishing of the derivatives of order $<m_\rho$ at the endpoint and a nonzero $m_\rho$-th derivative, rather than vanishing precisely for $j<m_\rho$ (\lbl{GJN0.8}, \lbl{GJNR8}); a proof step in \lbl{CS2} is done on differences before passing to classes. A wording correction to \lbl{TR0--TR6} (scalar trace values test zero location; the full trace operators recover local jets) could not be traced to a file with certainty. One claim was withdrawn in the programme on 24 September, together with the matrix file on which it rested (\note{41\_} §§1--2).
\item \emph{Small points found by the audit} in the programme's files, each harmless to the result concerned: typographical slips in the heat-metric blocks \lbl{HM10} and \lbl{HM19}; the expression $2j!$ in \lbl{NHI0}, \lbl{HSR0.2}, \lbl{HBW0.1} and \lbl{HBWR0.1} must be read as $2\cdot(j!)$; a stated file hash that differs by one character from the file's; and a corollary in one of the author's notes of July 2026 stating that the critical line, with the pairing $(s,1-\bar s)$, maps onto the whole Hopf fibre, whereas its image is the open half-circle (the line at $-\tfrac12$ gives the other half; \note{17\_} Proposition 17.3(c)).
\item \emph{The programme's Deligne reader} flagged a number of displays of Weil~II and repaired them. Its reading log records a comparison with the printed pages for §1, placing for example (1.3.10)(iv) and the twist in (1.6.14.3) in the original; its sections on §§2--3 make no claim about the printed original and treat the discrepancies there as defects of its transcription witnesses. The printed pages show that these, too, are in the original (Section~\ref{ssec:misprints}). The reader's mathematics stands.
\end{itemize}

\subsection{Corrections of printed sources}

The printed displays of Weil~II listed in Section~\ref{ssec:misprints}: six in §1 (two of them in (1.7.5)) and seven in §§2--3, among them the choice $\varepsilon_1+\varepsilon_2<\varepsilon$ in the proof of 2.1.7, which does not give Deligne's inequality, while $\varepsilon_1+2\varepsilon_2<\varepsilon$ does (Lemma~\ref{lem:concentration}). None affects a conclusion of Weil~II.

\subsection{Version~2 of this reader}

Version~1 of this reader (Zenodo 10.5281/zenodo.22970703) was refereed for errors and overstatement before publication. At the author's request it was then reviewed for understatement and missed generality as well as for overstatement, by a Claude instance that had written neither the reader nor the notes. The author of this reader verified its four corrections and five major findings against the sources, re-deriving each proof, with independent checks (Z1--Z6); a second independent instance verified the seventeen minor findings with checks of its own (\note{48\_}; \note{checks/zeta/zeta\_reader\_review/verify\_minors/}). Version~2 applies them; the status line of every changed statement says what changed. Version~2 then had its own referee pass (Appendix~\ref{app:verification}), whose two new mathematical items, the angle condition of Proposition~\ref{prop:NB} and composite moduli in Proposition~\ref{prop:SWzeros}, were re-derived by the author (checks Z7--Z8) and applied. \note{48\_} and the review use the numbering of version~1; version~1's Lemmas~2.8--2.9 and Corollary~2.10 are Lemmas~\ref{lem:shifted}--\ref{lem:convroot} and Corollary~2.11 here.

\emph{Corrections.}
\begin{itemize}
\item O1. Proposition~\ref{prop:avgweight} stated the rank $4m$ and the determinant $p^{2m}$ for every zero. At an on-line zero the four points $\rho,\bar\rho,1-\rho,1-\bar\rho$ are two, the rank is $2m$ and the determinant $p^m$. The conclusion, average weight $1$, is unchanged.
\item O2. Corollary~\ref{cor:eulerian} said that $\zeta(s,1+t)$ has an Euler product of Euler type only for $t\in\{0,-\tfrac12\}$. At $t=-\tfrac12$, $\zeta(s,\tfrac12)=2^s\prod_{p>2}(1-p^{-s})^{-1}$, so the Euler product is that of the normalised sheet $2^{-s}\zeta(s,\tfrac12)$; the corollary now speaks of $a^s\zeta(s,a)$.
\item O3. The source note \note{02\_} §4.2 concludes that both error terms of Corollary~B tend to $0$ once $\Delta\ge(1+\delta)\sqrt{2\log T_2}/\log2$ for a fixed $\delta>0$; this also needs $T_1/\Delta\to\infty$ (for example $T_1\ge4\Delta\sqrt{\log\Delta}$) and $O$-constants independent of $\Delta$. The bullet on Corollaries~B, B$_a$ in Section~\ref{sec:sheets} now says so.
\item O4. The outline of the proof of Proposition~\ref{prop:NB} ended with a step that gives $\Lambda=0$ only on $g_t\B$. The missing step, the density of $g_t\B$ in $\B$, is now proved, and the proposition has a complete proof.
\end{itemize}

\emph{Stronger statements.} Theorem~\ref{thm:jet} on every sheet $a>0$ with its whole-plane divisors; Lemma~\ref{lem:PTQ} for all bounded Hermitian forms compatible with the transfer for $n=2,3$; Proposition~\ref{prop:NB} with Carleman criteria in the half-plane and in the angle $|\arg w|<\pi/4$, which answer open question~5 for $t>(\log2)(\log3)/(4\pi)$ (the angle step is the version-2 referee's); Theorem~\ref{thm:formallog} and Corollary~\ref{cor:firstneg} at every element whose proper $M$-divisors factor uniquely; and Proposition~\ref{prop:abscissa} in every tensor degree, without positivity. These are register items S73--S77.

\emph{Smaller additions and omissions made good.} Proposition~\ref{prop:congruence} without Dirichlet's theorem and Example~\ref{ex:foureight} on the whole half-plane of absolute convergence; Proposition~\ref{prop:SWzeros} (zeros of $Z_{1/q}$ on both sides of the line for every $q$ with $\varphi(q)>2$, using Saias--Weingartner as stated in its abstract; S78); Lemmas~\ref{lem:shifted} and~\ref{lem:convroot} for complex data; Theorem~D's hypothesis in the verified range; Proposition~\ref{prop:spectral-L} and Lemma~\ref{lem:disjoint} for $D\le\operatorname{div}\Phi$; Lemma~\ref{lem:expsum} with polynomial factors; Lemma~\ref{lem:jetdecay} with a sharper factor and Corollary~\ref{cor:jetimage} with density in $\Lambda_m$; Theorem~\ref{thm:primary}(b$'$) and its form for $L(s,\chi)$; Proposition~\ref{prop:NBfail} with infinite codimension; the residue-duality constant $8.79$; Proposition~\ref{prop:blind} with one dilation; Proposition~\ref{prop:chiral}(c) unconditionally, with the $L(s,\chi)$ forms of Theorem~\ref{thm:mirror} and Proposition~\ref{prop:chiral} (S79); Proposition~\ref{prop:races} for every modulus; the separator for Dedekind zeta functions; the direction of monotonicity in Lemma~\ref{lem:defect}; Lemmas~\ref{lem:intpos} and~\ref{lem:powersums} in real-valued and exponential forms; Proposition~\ref{prop:clocks} for every integer $a\ge2$; and Lemmas~\ref{lem:escape} and~\ref{lem:keller} for real $a>1$ and for $C^1$ fields. The independent verification corrected the review's wording in six of these (Proposition~\ref{prop:SWzeros}, Theorem~D, Theorem~\ref{thm:primary}, Proposition~\ref{prop:chiral}, the Dedekind separator and Lemma~\ref{lem:escape}), and rejected one proposed item, for Proposition~\ref{prop:detect}, as not an understatement.

None of these is a new statement about the location of the zeros.


\subsection{Version~3 of this reader}

Version~2 was withdrawn on 27 September 2026, together with the audit notes, because of its framing. Version~3 keeps every theorem, proof, figure and citation of version~2 and corrects the framing; the only changes of content are the two decimal renderings and the three edge statuses listed at the end of this subsection. The changes of framing are of four kinds.

\emph{Scoped negative results instead of verdicts.} Where version~2 turned a scoped negative result into a general statement that a route does not decide the location of the zeros, version~3 states the scoped result and, where it is known, the input the construction would need. This concerns Section~\ref{sec:negative}, now titled \emph{Proved negative results, with their exact scope}, and in particular:
\begin{itemize}
\item the items N-Mirror, N-Logic, N-Pos1--N-Pos4, N-Unif1 and N-Unif3--N-Unif6, N-Prim, N-W1--N-W5, N-F1, N-Two3 and N-X;
\item the closing paragraph of Section~\ref{ssec:receivers};
\item the discussion after Proposition~\ref{prop:abscissa} and of Weil~II~1.5.1;
\item the status of Lemma~\ref{lem:escape};
\item the discussion after Proposition~\ref{prop:detect}.
\end{itemize}
The item N-Unif6 now states that the cohomology of the doubled source is that of Connes--Consani's base, and what this gives for the spectral realisation of the zeros, instead of a summary that the doubled source adds nothing.

\emph{Bridges.} Section~\ref{sec:bridges} is restructured. Each bridge now gives its typed map, what is established, what it suggests on each side and my assessment, and the introduction presents the bridges among the main outcomes. The mathematical content of each item of version~2's list is kept.

\emph{Front matter.} The abstract and introduction are rewritten: they state the main results as Theorems~I--V, explain what they mean, describe the bridges, and set out the four statuses. They also say how the research was carried out and where the withdrawn notes and the republished checks are.

\emph{Wording.} Quotation marks around the author's terms and around standard terms are removed. The word \emph{pure} in Section~\ref{sec:bridges} is used in its standard sense for Weil--Deligne representations. Session reports are cited as the programme's. Question~\ref{q:tau} and the other items of Section~\ref{sec:open} are stated as questions.

\emph{Changes of content.}
\begin{itemize}
\item Two decimal renderings were wrong. One has $(\log2)(\log3)/(4\pi)=0.0605982\ldots$ and $(\log2)(\log3)/(2\pi)=0.1211965\ldots$, which version~2 printed as $0.06060\ldots$ and $0.12120\ldots$ (in Section~\ref{sec:quotient}, after Proposition~\ref{prop:NB}, and in Question~\ref{q:sparse}); version~3 prints $0.060598\ldots$ and $0.121196\ldots$. The thresholds themselves, stated as the closed forms, are unchanged.
\item Three statuses in the table of Section~\ref{ssec:overlaps} are updated from the companion records: the edge from $S^6$ to the Erdős--Straus workbench (the orbits of the covers are verified at every level in the Erdős--Straus record), the edge from problem 817 to zeta (one of the ten interface notes is verified in the companion record), and the Collatz edge (verified in part in the companion record). Each keeps the status it had in version~2, marked as such.
\end{itemize}
Apart from the two decimal renderings and the three updated edge statuses, no statement changed its mathematical content, and none was removed.
